\documentclass[book.tex]{subfiles}
\begin{document}

\chapter{Euler's Method}

\label{chap:ode_euler}
\noindent\rule{11cm}{0.4pt} \\
\textbf{Prerequisites:} \autoref{chap:linear_systems} \\
\textbf{Difficulty Level:} ** \\
\noindent\rule{11cm}{0.4pt}
\vspace{5mm}

In this chapter, we learn how to implement the Euler's method for solving a single ordinary differential equation. We learn about the meanings of local and global truncation errors and their relationships to step sizes for one-step methods for solving ODEs. We then solve the Cahn-Hilliard model using Euler's method and end by explaining the use of adaptive step-sizing for solving ODEs.

\section{First Order ODEs}

First-order ODEs form the basis for solving governing equations of engineering systems. Let us first consider different approaches for solving first-order ODEs. %As discussed in the previous lecture, 
One-step methods and multi-step methods are two widely used methods for solving these first order ODEs. 

\textit{One-Step methods} attempt to trace out the trajectory of the solution into the future by extrapolating from an old value $y_i$ to a new value $y_{i+1}$ over a distance $h$, given the value of the increment function (first-derivative of the function) at a single point $i$. They are also referred to as Runge-Kutta methods after the two applied mathematicians who first discussed them in the early 1900s. \textit{Multi-Step methods} use information from several previous points as the basis for extrapolating to a new value. We focus on studying solution methods for first order ODEs since higher order ODEs can be reduced to a system of first order ODEs.  %We will be discussing one-step methods in particular for the next two lectures. 

Runge-Kutta methods are a family of iterative methods, usually employed for solving first order ODEs. Among these methods, the simplest approach to solve a first order ODE is to use the differential equation to estimate the slope in the form of the first derivative at $t_i$. In other words, the slope at the beginning of the interval is taken as an approximation of the average slope over the whole interval. This approach is called the \textit{Euler’s method}. It is the most basic method for numerical integration of ordinary differential equations and is the simplest Runge–Kutta method.
%\end{document}
\section{Euler's Method}

Let's look at the first order ODEs of the general form $\frac{dy}{dt} = f(t,y)$. The first derivative of the function provides a direct estimate of the slope at the initial value $(t_i,y_i)$. This is given as 
\begin{align}\phi = \frac{dy}{dt} = f(t_i,y_i),\end{align}
where $f(t_i,y_i)$ is the differential equation evaluated at $t_i$ and $y_i$. The estimate from Euler's method can be obtained by moving over by a step size along this slope estimate from the initial value. The estimate at the next step is given as 
\begin{align}
y_{i+1} = y_{i} + f(t_i,y_i)h
\end{align}
\begin{marginfigure}[5\baselineskip] 
	\centering
	% \includegraphics[width = 2.5in]{figures/part1b/euler_method/Euler.png}
 \includegraphics[width = 2.5in]{figures/part1b/euler_method/Euler_0.png}
	\centering
	\caption{Euler's Method}
	\label{fig:1}
\end{marginfigure}

This formula is referred to as \textit{Euler’s method} (or the Euler-Cauchy or point-slope method). A new value of y is predicted using the slope (equal to the first derivative at the original value of t) to extrapolate linearly over the step size $h$. It can be observed that the Euler's method primarily draws its motivation from the Taylor's series by utilizing the first two-terms of the function's approximation about the initial value.


\section{The Cahn-Hilliard Equation}
Let's consider the Cahn-Hilliard Equation which is an example of a non-linear differential equation describing the process of phase separation, by which the two components of a binary fluid spontaneously separate and form domains which are pure in each component. The Cahn-Hilliard Equation is given as
\begin{align}
\dfrac{\partial c}{\partial t} = D \nabla^2 (c^3 - c - \gamma \nabla^2c )
\vspace{-6pt}
\end{align}
Let's try to solve the Cahn-Hilliard equation given above using the Euler's method from $t=0$ to $t=3$ with initial conditions as the normalized intensity values of an image. Suppose time step $dt$ of $0.002$ s, a diffusion coefficient $D$ of $20$ and $\gamma$ of $0.5$.

The Cahn-Hilliard equation can be re-written as 
\begin{align}
\dfrac{\partial c}{\partial t} = D (\nabla^2c^3 - \nabla^2 c - \nabla^2 \gamma \nabla^2 c)
\end{align}

We assume for modeling purposes that we can model $c$ by a $N\times N$ sized grid.

\begin{lstlisting}[language=python]
def cahn_hilliard_eulers(C : $\mathbb{R}^{N\times N}$, $\gamma: \mathbb{R}$, D : $\mathbb{R}$, dt: $\mathbb{R}$, ($t_0$, $t_e$): $\mathbb{R}^2$, tol: $\mathbb{R}$) $\to$ $\mathbb{R}^{N\times N}$:
  out = zeros(N, N)
  $t$ = $t_0$ 
  while $t$ < $t_e$:
    dc$^2$ = $\nabla^2$(C, dx) + $\nabla^2$(C, dy)
    dc$^3$ = $\nabla^2$(C$^3$,dx) + $\nabla^2$(C$^3$,dy)
    $\frac{\textrm{d}\gamma}{\textrm{dc}}$ = $\nabla^2$(gamma*dc$^2$, dx) + $\nabla^2$(gamma*dc$^2$, dy)
    $\frac{\textrm{dc}}{\textrm{dt}}$ = D*(dc$^3$ - dc$^2$ - $\frac{\textrm{d}\gamma}{\textrm{dc}}$)
    C = C + $\frac{\textrm{d}\gamma}{\textrm{dc}}$*dt
    if $|\|\textrm{out}\|-\|C\|)|$ < tol:
      break
    out = C
    time = time + dt
  return out
\end{lstlisting}
We use the notation $\nabla^2$ above, but in practice we need to use the central difference to approximate the second-order derivative at a grid element. %The function \textit{bound\_conc}() ensures that the concentration lies within [-1,1].
%\end{document}
\section{Error Analysis for Euler's Method}
The numerical solution of ODEs involves two types of error
\begin{enumerate}
	\item \textit{Truncation} or Discretization Errors - These are caused by the nature of the techniques employed to approximate values of y.
	\item \textit{Roundoff} Errors - These are caused by the limited numbers of significant digits that can be retained by a computer.
\end{enumerate}
The truncation errors are composed of two parts. The first is a \textit{local truncation error} that results from an application of the method in question over a single step. The second is a \textit{propagated truncation error} that results from the approximations produced during the previous steps. The sum of the two is the total error. It is referred to as the \textit{global truncation error}.

Insight into the magnitude and properties of the truncation error can be gained by deriving
Euler’s method directly from the Taylor series expansion. To do this, realize that the
differential equation being integrated will be of the general form of $\frac{dy}{dt} = f(t,y)$ with an initial value of ($t_i,y_i$). If the solution, i.e. the function describing the behavior of y has continuous derivatives, it can be represented by a Taylor series expansion about a starting value ($t_i ,y_i$), as in
\begin{align}
y_{i+1} = y_{i} + y_{i}'h + \frac{y_i''}{2!}h^2 + \dots + \frac{y_i^{n}}{n!}h^n + R_n
\end{align}

where $h = t_{i+1} - t_{i}$ and $R_n = $ the remainder term defined as
\begin{align}
R_n = \frac{y^{n+1}(\xi)}{(n+1)!}h^{n+1}
\end{align}
where $\xi$ lies somewhere in the interval from $t_i$ to $t_{i+1}$. The above Taylor Series expansion can also be re-written as
\begin{align}
y_{i+1} = y_{i} + f(t_i,y_i)h + \frac{f'(t_i,y_i)}{2!}h^2 + \dots + \frac{f^{(n-1)}(t_i,y_i)}{n!}h^n + O(h^{n+1})
\end{align}

where $O(h^{n+1})$ specifies that the local truncation error is proportional to the step size raised to the $(n + 1)$th power. From the above equation, it can be seen that the Euler's method covers only the first two terms of the Taylor's expansion. Thus, the truncation error in Euler’s method is attributable to the remaining terms in the Taylor series expansion that were not included and is given as
\begin{align}
E_t = \frac{f'(t_i,y_i)}{2!}h^2 + \dots + \frac{f^{(n-1)}(t_i,y_i)}{n!}h^n + O(h^{n+1})
\end{align}

where $E_t$ is the \textit{true local truncation error}. For sufficiently small $h$, the higher-order terms are usually negligible, and the result is often represented as 
\begin{align}
E_a = \frac{f'(t_i,y_i)}{2!}h^2 = O(h^2)
\end{align}

where $E_a$ is the \textit{approximate local truncation error}. 

It can also be shown that the global truncation error is $O(h)$, i.e. it is proportional to the step size. These observations lead to some useful conclusions:
\begin{enumerate}
	\item The global error can be reduced by decreasing the step size.
	\item The method will provide error-free predictions if the underlying function is linear, because for a straight line the second derivative would be zero.
\end{enumerate}
This latter conclusion makes intuitive sense because Euler’s method uses straight-line segments to approximate the solution. Hence, Euler’s method is referred to as a \textit{first-order method}.

\section{Stability of Euler's Method}

The stability of a solution method is another important consideration that must be considered when solving ODEs. A numerical solution is said to be unstable, if errors grow exponentially for a problem for which there is a bounded solution. The stability of a particular application can depend on three factors: the differential equation, the numerical method, and the step size.

Insight into the step size required for stability can be examined by studying the following simple ODE:
\begin{align}
\frac{dy}{dt} = -ay
\end{align}

If $y(0)=y_0$, the true solution can be determined using calculus as
	
\begin{align}
y = y_0e^{-at}
\end{align}

The above solution starts at $y_0$ and asymptotically approaches zero. Solving the same problem using Euler's method gives the following result,
	
\begin{align}
y_{i+1} &= y_i + \frac{dy_i}{dt}\\
h &= y_i - ay_ih = y_i(1-ah)
\end{align}

The parenthetical quantity $1-ah$ is called an \textit{amplification factor}. If its absolute value is greater than unity, the solution will grow in an unbounded fashion. So clearly, the stability depends on the step size \textit{h}. Based on this analysis,
Euler’s method is said to be \textit{conditionally stable}.

Note that there are certain ODEs where errors always grow regardless of the method.
Such ODEs are called \textit{ill-conditioned}. Inaccuracy and instability are often confused. This is usually because both represent situations where the numerical solution breaks down or both are affected by step size.
%\end{document}
\section{Adaptive Step-Sizing}

We have seen how to solve ODEs numerically using Euler's method using
%and Runge-Kutta's methods 
a fixed step size. However, it is noticed that the local trunctation error is proportional to the step-size chosen. By changing this step-size dynamically at every step, an upper bound on the local truncation error can be imposed with minimum number of steps. \textit{Adaptive step-sizing} methods make use of this to come up with a new step-size at every step.

The local truncation error for Euler's method is given by $E_a = \frac{f'(t_i,y_i)}{2!}h^2$. The step-size is simply obtained by establishing an upper bound on this truncation error ($e_{max}$).
	\begin{align}
	E_a = \frac{f'(t_i,y_i)}{2!}h^2 <= e_{max} \implies h <= \sqrt{\dfrac{2e_{max}}{f'(t_i,y_i)}}
	\vspace{-6pt}
	\end{align}
	
This condition is used for limiting the local truncation error while achieving greater performance on segments where the derivative of $f$ is small. 

%For second order Runge-Kutta method, the local truncation error is given by $E_r = \frac{f''(t_i,y_i)}{3!}h^3$, The criterion for adaptive step size can be similarly obtained here as
%	\vspace{-6pt}
%	\begin{equation*}
%	E_r = \frac{f''(t_i,y_i)}{3!}h^3 <= e_{max} \implies h <= \left(\dfrac{6e_{max}}{f''(t_i,y_i)}\right)^{1/3}
%	\vspace{-6pt}
%	\end{equation*}
	
As an example, let's try to estimate $y$ from the ODE 
\begin{align}
\frac{dy}{dt} = t-\log(1+t)
\end{align}
using adaptive step-sizing and Euler's method.

\begin{marginfigure} 
	\centering
	% \includegraphics[width = 2.5in]{figures/part1b/euler_method/Euler.png}
 \includegraphics[width = 2.5in]{figures/part1b/euler_method/Euler_0.png}
	\centering
	\caption{Step size for Euler's Method}
	\label{fig:3}
\end{marginfigure}
\begin{lstlisting}[language=python]
def adaptive_euler(f: $\mathbb{R} \to \mathbb{R}$, $(t_0, t_e)$: $\mathbb{R}^2$, tol: $\mathbb{R}$) $\to \mathbb{R}[]$:
  t_adp, y, dt: $\mathbb{R}[]$
  t_adp[0] = $t_0$, i = 0
  y[0] = 0
  dt[0] = sqrt(2*tol/$\nabla$f(t_adp[0]))
  while (t_adp[i] < $t_e$):
    y[i+1] = y[i] + f(t_adp[i])*dt[i]
    dt[i+1] = sqrt(2*tol/$\nabla$f(t_adp[i]))
    t_adp[i+1] = t_adp[i] + dt[i+1]
    i = i + 1
  y
\end{lstlisting}
It can be observed that initially the step-size was high, which gradually decreases because of the increasing value of $f'$.
\end{document}